\section{Integral monodromy: the kernel calculation}\label{app:cdp-monodromy}
We retain the complete kernel calculation for Theorem~\ref{fam:st:F.1.1}.
\begin{proof}
\textit{Step 1: the integral stabilizer.}
Since $Ge=e$ and $G$ preserves $\Omega$, the identity
$\Omega(Gv,e)=\Omega(v,e)$ forces the first row of $G$ to be
$(1,0,0,0)$. The last column is $e$, so write
\[
 G=\begin{pmatrix}1&0&0\\v&B&0\\t&w&1\end{pmatrix}.
\]
Preservation of $\Omega$ is equivalent to
$B^tJB=J$ and $6w+v^tJB=0$. The first equality gives
$B\in\operatorname{SL}_2(\mathbb Z)$. The second and integrality of $w$
give $v^tJB\in6\mathbb Z^2$. Because $JB$ is unimodular,
$v=6x$ for a unique integral vector $x$. This proves the asserted
parametrization, including its converse. Block multiplication gives
\begin{equation}\label{cdp:eq:parabolic-law}
 g(B,x,t)g(C,y,s)=g(BC,x+By,t+s-6x^tJBy).
\end{equation}

\textit{Step 2: the cyclic character of the middle block.}
Put
\[
 S=\begin{pmatrix}0&-1\\1&0\end{pmatrix},\qquad
 a=\begin{pmatrix}0&1\\-1&-1\end{pmatrix},\qquad
 T=Sa=\begin{pmatrix}1&1\\0&1\end{pmatrix}.
\]
The Euclidean algorithm on the first column of an integral
unimodular matrix shows that $S,T$, hence $a,S$, generate
$\operatorname{SL}_2(\mathbb Z)$. Their images in
$\operatorname{PSL}_2(\mathbb Z)$ have orders three and two.
For the fractional-linear action on $\mathbb R\cup\{\infty\}$,
$a$ sends $(0,\infty)$ into $(-1,0)$, $a^2$ sends
$(0,\infty)$ into $(-\infty,-1)$, and $S$ sends
$(-\infty,0)$ into $(0,\infty)$. The reduced-word ping-pong argument
therefore identifies the subgroup they generate in the projective
group with the free product of their cyclic subgroups: a reduced
alternating word, after conjugation to start and end in the
order-three factor, sends the positive interval into the negative
interval and cannot be the identity. Such a conjugation is possible
for every nontrivial cyclically reduced alternating word because
the order-three factor has two nonidentity elements; a word in one
factor is already nonidentity. Lifting to the integral group gives
\[
 \operatorname{SL}_2(\mathbb Z)
 =\langle a,S\mid a^3=S^4=1,\ [a,S^2]=1\rangle.
\]
Indeed, the presented group maps surjectively onto the integral
group; quotienting its central element $S^2$ gives the just-proved
free product, and $S^2$ maps to $-I$, so the map has trivial kernel.
The assignments $\alpha(a)=4$ and $\alpha(S)=9$ respect these
relations modulo $12$ and give $\alpha(T)=1$. They define the
required homomorphism, which is unique because $S,T$ generate.

\textit{Step 3: the character of the stabilizer.}
Modulo two, $Q$ vanishes exactly at the origin of $(\mathbb Z/2)^2$.
Consequently $Q(Bx)=Q(x)$ modulo two for every integral unimodular
$B$. Also
\[
 Q(x+y)-Q(x)-Q(y)=x_1y_2+x_2y_1=x^tJy\pmod2.
\]
Together with \eqref{cdp:eq:parabolic-law} these identities show that
\[
 \chi(g(B,x,t))=t+6Q(x)-\alpha(B)\pmod {12}
\]
is a homomorphism. It is surjective because
$\chi(g(I,0,1))=1$. Direct substitution gives
\[
 A_1=g(a,(1,-1)^t,-2),\qquad A_2=g(S,(0,-1)^t,3),
 \qquad \chi(A_1)=\chi(A_2)=0.
\]
Thus $\langle A_1,A_2\rangle\subseteq\ker\chi$.

\textit{Step 4: equality with the kernel.}
With $[G,H]=GHG^{-1}H^{-1}$, multiplication in
\eqref{cdp:eq:parabolic-law} gives
\[
 \begin{aligned}
 K&=[A_1,A_2^2]=g(I,(0,-1)^t,6),\\
 U&=A_2KA_2^{-1}=g(I,(1,0)^t,-6),\\
 V&=K^{-1}=g(I,(0,1)^t,-6),\\
 Z&=I+12e\gamma=g(I,0,12),\qquad [U,V]=Z^{-1}.
 \end{aligned}
\]
In particular $U,V,Z$ belong to $\langle A_1,A_2\rangle$. For all
integers $r,s,n$, the multiplication law gives
\[
 U^rV^sZ^n=g(I,(r,s)^t,-6(r+s+rs)+12n).
\]
These are precisely the elements of $\ker\chi$ whose middle block
is $I$. The middle blocks $a,S$ of $A_1,A_2$ generate
$\operatorname{SL}_2(\mathbb Z)$. Given any $G\in\ker\chi$, choose
$H\in\langle A_1,A_2\rangle$ with the same middle block as $G$.
Then $H^{-1}G$ belongs to the already computed kernel over $I$ and
hence to $\langle A_1,A_2\rangle$. Therefore
$\ker\chi=\langle A_1,A_2\rangle$, and surjectivity of $\chi$ proves
the index formula.
\end{proof}
